\renewcommand{\baselinestretch}{1.75}\selectfont

\chapter{کلیات و مبانی تشخیص توضیح‌پذیر سلامت روان}

\section{مقدمه و بیان مسئله}

سلامت روان بخشی جدایی‌ناپذیر از سلامت عمومی است و اختلال در آن می‌تواند کیفیت زندگی، روابط اجتماعی، تحصیل و اشتغال فرد را تحت تأثیر قرار دهد. در بسیاری از جوامع، نیاز به خدمات سلامت روان بیش از ظرفیت پاسخ‌گویی نظام‌های درمانی است \cite{who2022mentalhealth}. این فاصله، ضرورت توسعه ابزارهایی را نشان می‌دهد که شناسایی اولیه و بررسی شواهد را برای متخصصان تسهیل کنند، بدون آن‌که جایگزین ارزیابی بالینی شوند.

تشخیص اختلالات روانی بر پایه مصاحبه بالینی، مشاهده، سابقه فرد و معیارهای تشخیصی انجام می‌شود و به مدت و شدت نشانه‌ها و اثر آن‌ها بر زندگی روزمره وابسته است \cite{kupfer2013dsm5}. در کنار این منابع، نوشته‌های منتشرشده در فضای مجازی حاوی گزارش‌هایی از تجربه‌ها و حالات هیجانی کاربران هستند و می‌توان از آن‌ها برای مطالعه الگوهای مرتبط با افسردگی، اضطراب و تنیدگی استفاده کرد \cite{shen2017depression,sarkar2023predicting}. با این حال، چنین متنی معادل پرونده بالینی نیست و نتیجه حاصل از آن باید یک برآورد محاسباتی تلقی شود.

روش‌های طبقه‌بندی معمولاً یک برچسب مانند وجود یا نبود افسردگی تولید می‌کنند، اما نشان نمی‌دهند کدام شواهد متنی در تصمیم نقش داشته‌اند. این محدودیت در سلامت روان مهم است؛ زیرا باید بتوان برداشت نادرست از کنایه، نادیده گرفتن بافت یا تعمیم از یک نشانه منفرد را تشخیص داد. از این رو، توضیح‌پذیری\footnote{\lr{Explainability}} بخشی از ارزیابی مسئولانه سامانه است، نه صرفاً یک قابلیت جانبی \cite{joyce2023explainable}.

مدل‌های زبانی بزرگ\footnote{\lr{Large Language Models (LLMs)}} می‌توانند همراه با برچسب، توضیحی متنی نیز تولید کنند. با وجود این، روان بودن توضیح درستی آن را تضمین نمی‌کند و مدل ممکن است نشانه یا رابطه‌ای را بیان کند که در ورودی وجود ندارد. این مسئله که توهم اطلاعاتی\footnote{\lr{Factual Hallucination}} نامیده می‌شود، در حوزه سلامت روان می‌تواند به نتیجه‌گیری قطعی بر پایه شواهد ناکافی بینجامد. به همین دلیل، شفافیت، ارزیابی دقیق و نظارت انسانی در کاربردهای سلامت ضروری است \cite{who2021aiethics,who2023safeai}.

بازیابی افزوده\footnote{\lr{Retrieval-Augmented Generation (RAG)}} با قراردادن اطلاعات بیرونی در اختیار مدل، راهی برای کاهش وابستگی به دانش درونی آن فراهم می‌کند \cite{lewis2020retrieval}. اگر این اطلاعات در قالب گراف دانش\footnote{\lr{Knowledge Graph (KG)}} سازمان‌دهی شوند، روابط میان مفاهیم نیز در بازیابی و استدلال قابل استفاده خواهند بود. مسئله این پژوهش، بازپیاده‌سازی و انطباق چند روش گراف‌محور و مقایسه آن‌ها در یک وظیفه مشترک سلامت روان است. روش‌ها از نظر درستی پیش‌بینی، اتکاپذیری توضیح، زمان پاسخ و هزینه متنی بررسی می‌شوند؛ هدف نیز ارائه یک روش گرافی تازه نیست، بلکه روشن‌کردن مزایا و محدودیت‌های الگوهای موجود است.

\section{تشخیص توضیح‌پذیر سلامت روان از داده‌های متنی}

در این پایان‌نامه، تشخیص از متن به معنای نگاشت یک نوشته خوداظهاری به برچسبی از پیش تعریف‌شده است. ورودی می‌تواند یک پیام کوتاه یا نوشته‌ای منتشرشده در شبکه اجتماعی باشد و خروجی شامل برچسب وضعیت روانی یا عامل مورد بررسی و توضیح شواهد مؤثر در آن است. منظور از تشخیص توضیح‌پذیر، صدور حکم بالینی نیست؛ سامانه باید تصمیم محاسباتی خود را به شکلی ارائه دهد که قابل بازبینی باشد.

زبان این نوشته‌ها اغلب غیررسمی، استعاری یا غیرمستقیم است و معنای یک عبارت می‌تواند به جمله‌های پیشین و پسین وابسته باشد. افزون بر این، نشانه‌هایی مانند اختلال خواب، خستگی و کاهش تمرکز به یک اختلال خاص محدود نیستند. بنابراین، مدل باید میان شواهد صریح، برداشت احتمالی و اطلاعاتی که در متن وجود ندارد تمایز بگذارد.

مجموعه‌داده \lr{IMHI}\footnote{\lr{Interpretable Mental Health Instruction}} داده‌های چند مجموعه پیشین را در قالب تعامل‌های دستور و پاسخ گردآوری کرده است \cite{yang2023towards,yang2024mentallama}. وظایف آن شامل تشخیص دودویی، تشخیص چندکلاسه، شناسایی علت یا عامل زمینه‌ای و تشخیص عوامل خطر یا ابعاد بهزیستی است. در این پژوهش، از هر گروه یک تنظیم نماینده انتخاب می‌شود تا روش‌ها در صورت‌های متفاوت مسئله و در عین حال در دامنه‌ای قابل کنترل مقایسه شوند.

توضیح مطلوب باید به متن ورودی وفادار باشد و از افزودن نشانه، رویداد یا رابطه‌ای که در متن دیده نمی‌شود پرهیز کند. ممکن است برچسب نهایی درست باشد، اما مسیر استدلال شامل ادعایی نادرست باشد؛ از این رو، ارزیابی برچسب و توضیح باید جداگانه انجام شود. همچنین، خطا تنها به کاهش یک معیار عددی محدود نیست: مثبت کاذب می‌تواند متن عادی را نشانه اختلال تلقی کند و منفی کاذب ممکن است شواهد مهم را نادیده بگیرد.

گراف دانش دامنه‌ای در این پژوهش منبعی صریح برای نگهداری مفاهیمی مانند اختلال، نشانه و عامل خطر و روابط میان آن‌هاست. روش گراف‌محور می‌تواند اطلاعات مرتبط را بازیابی یا مسیرهای میان مفاهیم را دنبال کند، اما وجود گراف به‌تنهایی درستی پاسخ را تضمین نمی‌کند؛ کیفیت بازیابی، پوشش گراف و نحوه استفاده مدل از شواهد بر نتیجه اثر دارند. سامانه نیز صرفاً یک ابزار تصمیم‌یار است و نتایج آن در محدوده داده‌ها و تنظیمات آزمایشی تفسیر می‌شود. حریم خصوصی، رضایت، سوگیری داده و تفاوت‌های زبانی و فرهنگی نیز باید در استفاده از داده‌های سلامت روان در نظر گرفته شوند \cite{who2021aiethics}.

\section{مدل‌های زبانی بزرگ، توضیح‌پذیری و توهم اطلاعاتی}

مدل زبانی، توزیع احتمال دنباله‌های زبانی را می‌آموزد و در زمان تولید، هر جزء متن را با توجه به اجزای پیشین پیش‌بینی می‌کند. مدل‌های زبانی بزرگ این سازوکار را با تعداد زیادی پارامتر و بر پیکره‌های گسترده به کار می‌گیرند؛ در نتیجه، می‌توانند متن ورودی را خلاصه کنند، به پرسش پاسخ دهند و برای یک تصمیم توضیح زبانی بسازند. در این پایان‌نامه، اهمیت این مدل‌ها در توانایی انجام هم‌زمان تشخیص و تولید توضیح است، نه در استفاده از آن‌ها به عنوان منبع مستقل دانش بالینی.

دانش مدل در پارامترهای آن ذخیره شده و دسترسی مستقیم به منشأ هر گزاره آسان نیست. پاسخ نیز به صورت پرسش، اطلاعات موجود در بافت ورودی و شیوه تنظیم مدل وابسته است. این ویژگی انعطاف‌پذیری زیادی ایجاد می‌کند، اما سبب می‌شود دو پاسخ زبانی مشابه از نظر پشتوانه دانشی یکسان نباشند. مدل ممکن است از یک عبارت مبهم نتیجه‌ای فراتر از متن بگیرد یا دانشی کلی را بدون مشخص‌کردن منبع وارد پاسخ کند.

در تشخیص توضیح‌پذیر، توضیح باید نشان دهد چه شواهدی به برچسب منتهی شده‌اند و این شواهد چگونه با نتیجه ارتباط دارند. توضیح روان یا متقاعدکننده لزوماً توضیح وفادار نیست؛ ممکن است پس از تولید برچسب ساخته شده باشد و روند واقعی تصمیم را منعکس نکند. بنابراین، کیفیت نگارش و وفاداری دو معیار متفاوت‌اند. اولی به وضوح و انسجام پاسخ مربوط است و دومی بررسی می‌کند که ادعاهای توضیح با متن ورودی و دانش معتبر پشتیبانی می‌شوند یا نه \cite{joyce2023explainable}.

توهم اطلاعاتی زمانی رخ می‌دهد که مدل گزاره‌ای فاقد پشتوانه در ورودی یا منبع معتبر تولید کند. در مسئله این پژوهش، این خطا می‌تواند به شکل افزودن نشانه‌ای گزارش‌نشده، نسبت‌دادن علت به یک وضعیت، یا بیان رابطه‌ای نادرست میان نشانه و اختلال ظاهر شود. چنین خطایی حتی با وجود برچسب درست نیز اعتبار توضیح را کاهش می‌دهد. حساسیت داده‌های سلامت روان نیز سبب می‌شود ادعاهای قطعی و بدون شاهد با دقت بیشتری کنترل شوند \cite{who2023safeai}.

برای کاهش این مسئله، لازم است شواهد مورد استفاده مدل تا حد امکان صریح و قابل ردیابی باشند. قرار دادن اطلاعات بازیابی‌شده در بافت ورودی، امکان مقایسه ادعاهای پاسخ با یک منبع بیرونی را فراهم می‌کند؛ هرچند انتخاب شواهد نامرتبط یا استفاده نادرست از آن‌ها همچنان می‌تواند پاسخ را منحرف کند. بخش بعدی سازوکار بازیابی افزوده و نقش آن در پیوند مدل زبانی با منابع دانشی را بررسی می‌کند.

\section{بازیابی افزوده: تعریف، سازوکار و مزایا}

بازیابی افزوده روشی برای پیوند مدل زبانی با یک منبع دانشی بیرونی است. در این روش، مدل پاسخ را تنها بر پایه دانشی که هنگام آموزش در پارامترهای آن ذخیره شده است تولید نمی‌کند؛ ابتدا اطلاعات مرتبط با پرسش از یک مجموعه اسناد بازیابی می‌شود و سپس پرسش همراه با اطلاعات بازیابی‌شده در اختیار مدل قرار می‌گیرد. به این ترتیب، مرحله بازیابی مسئول یافتن شواهد و مرحله تولید مسئول ترکیب شواهد و ساخت پاسخ است \cite{lewis2020retrieval}.

در آماده‌سازی منبع دانش، اسناد معمولاً به قطعه‌های کوچک‌تر تقسیم می‌شوند. این فرایند که قطعه‌بندی\footnote{\lr{Chunking}} نام دارد، امکان جست‌وجوی بخش مرتبط یک سند طولانی را فراهم می‌کند. سپس هر قطعه با یک بردار نهفته\footnote{\lr{Embedding}} نمایش داده و همراه با متن و اطلاعاتی مانند شناسه سند در یک نمایه ذخیره می‌شود. در زمان پرسش، پرسش نیز به همان فضای برداری نگاشت می‌شود و قطعه‌هایی که بیشترین شباهت معنایی را دارند به عنوان نامزد انتخاب می‌شوند.

قطعه‌های برگزیده به بافت ورودی مدل افزوده می‌شوند و دستور تولید پاسخ را تشکیل می‌دهند. مدل باید با اتکا به این بافت، اطلاعات پراکنده را کنار هم قرار دهد و پاسخ نهایی را بسازد. تعداد قطعه‌های بازیابی‌شده باید محدود باشد؛ زیرا افزودن شواهد کم، احتمال حذف اطلاعات لازم را افزایش می‌دهد و افزودن شواهد زیاد می‌تواند اطلاعات نامرتبط وارد پاسخ کند یا از ظرفیت بافت مدل فراتر رود. بنابراین، کیفیت پاسخ تنها به توان مدل زبانی وابسته نیست و انتخاب درست شواهد نقش تعیین‌کننده‌ای دارد.

مزیت اصلی این رویکرد، جداسازی نسبی دانش از پارامترهای مدل است. منبع بیرونی را می‌توان بدون آموزش دوباره مدل اصلاح یا به‌روزرسانی کرد و در صورت نگهداری شناسه اسناد، منشأ شواهد را نیز بررسی نمود. همچنین، محدودکردن پاسخ به اطلاعات بازیابی‌شده می‌تواند از تولید برخی ادعاهای بدون پشتوانه بکاهد. این مزیت در حوزه سلامت روان مهم است؛ زیرا باید بتوان مشخص کرد هر بخش از توضیح به کدام داده متنی یا مفهوم دامنه‌ای متکی است.

با این حال، بازیابی افزوده تضمین نمی‌کند که پاسخ همواره درست باشد. اگر قطعه مرتبط بازیابی نشود، اگر متن بازیابی‌شده مبهم یا ناسازگار باشد، یا اگر مدل از شواهد به‌درستی استفاده نکند، خطا همچنان رخ می‌دهد. از این رو، بازیابی صرفاً مرحله‌ای پیش از تولید نیست، بلکه بخشی از زنجیره استدلال و یکی از مؤلفه‌های اصلی ارزیابی سامانه است. در ادامه، محدودیت‌های بازیابی صرفاً متنی و دلیل استفاده از ساختار گرافی بررسی می‌شود.

\section{محدودیت‌های بازیابی متنی و ضرورت ساختار گرافی}

بازیابی متنی معمولاً هر قطعه را مستقل از سایر قطعه‌ها رتبه‌بندی می‌کند. این سازوکار برای پرسش‌هایی که پاسخ آن‌ها در یک بخش مشخص از متن قرار دارد مناسب است، اما در پرسش‌های پیچیده‌تر ممکن است شواهد لازم میان چند سند پراکنده باشند. در چنین حالتی، شباهت معنایی میان پرسش و هر قطعه به‌تنهایی برای کشف ارتباط میان شواهد کافی نیست و قطعه‌ای که واژگان مشابهی ندارد، با وجود نقش آن در پاسخ، کنار گذاشته می‌شود.

قطعه‌بندی نیز با یک مبادله همراه است. قطعه‌های بزرگ بافت بیشتری حفظ می‌کنند، اما بازیابی آن‌ها اطلاعات نامرتبط بیشتری وارد ورودی مدل می‌کند. قطعه‌های کوچک دقیق‌ترند، ولی ممکن است رابطه میان جمله‌ها، بخش‌های یک سند یا چند سند را از بین ببرند. این مسئله در مجموعه‌های بزرگ که اطلاعات مربوط به یک مفهوم در نقاط مختلف توزیع شده است شدیدتر می‌شود و بازیابی متن را در دنبال‌کردن زنجیره‌های چندمرحله‌ای محدود می‌کند \cite{ma2025tog2,zhuang2026linearrag}.

بسیاری از پرسش‌ها به استدلال چندگامی\footnote{\lr{Multi-hop Reasoning}} نیاز دارند؛ یعنی رسیدن به پاسخ مستلزم عبور از چند رابطه میان اشخاص، رویدادها یا مفاهیم است. یک فهرست رتبه‌بندی‌شده از قطعه‌ها این روابط را به‌صورت صریح نمایش نمی‌دهد و مدل باید آن‌ها را پس از بازیابی بازسازی کند. این وابستگی احتمال نادیده‌ماندن یک حلقه از زنجیره یا ایجاد ارتباطی بدون پشتوانه را افزایش می‌دهد. پژوهش‌های گراف‌محور نیز ضعف بازیابی برداری را در بازنمایی ارتباط‌های ساختاری و حافظه تداعی‌گر گزارش کرده‌اند \cite{gutierrez2025hipporag2}.

ساختار گرافی این امکان را فراهم می‌کند که مفاهیم به صورت گره و ارتباط‌های میان آن‌ها به صورت یال نمایش داده شوند. پس از شناسایی یک مفهوم در پرسش، بازیابی می‌تواند به مفاهیم همسایه گسترش یابد و به جای چند قطعه مستقل، یک زیرگراف مرتبط در اختیار مدل قرار دهد. در نتیجه، هم محتوای متنی و هم مسیر ارتباط میان شواهد قابل استفاده خواهد بود. برای نمونه، \lr{G-Retriever} مسئله انتخاب بخش مرتبط یک گراف متنی را به صورت یافتن یک زیرگراف متصل صورت‌بندی می‌کند تا حجم ورودی کاهش یابد و ساختار لازم برای پاسخ حفظ شود \cite{he2024gretriever}.

با این همه، گراف راه‌حلی بدون خطا نیست. گراف دانش ممکن است ناقص باشد و استخراج خودکار موجودیت‌ها و روابط می‌تواند یال‌های نادرست یا ناسازگار ایجاد کند. بازیابی گرافی نامناسب نیز ممکن است نویز را از یک گره به بخش‌های دیگر گسترش دهد. بنابراین، ضرورت استفاده از گراف به معنای کنارگذاشتن متن نیست؛ هدف، افزودن سازمان ساختاری به شواهد متنی و فراهم‌کردن مسیرهای قابل ردیابی است. تفاوت روش‌های موجود عمدتاً به نحوه ساخت گراف، انتخاب زیرگراف و چگونگی ادغام آن با مدل زبانی بازمی‌گردد که در بخش بعد بررسی می‌شود.

\section{گراف دانش، الگوهای ادغام و جهت‌گیری پژوهش}

گراف دانش، اطلاعات را به صورت مجموعه‌ای از موجودیت‌ها و روابط میان آن‌ها نمایش می‌دهد. هر گزاره را می‌توان با یک سه‌تایی شامل موجودیت مبدأ، نوع رابطه و موجودیت مقصد بیان کرد. برای نمونه، ارتباط میان یک نشانه، عامل خطر و اختلال می‌تواند به صورت چند سه‌تایی مرتبط ذخیره شود. برخلاف متن پیوسته، این بازنمایی امکان جست‌وجوی مستقیم همسایه‌های یک مفهوم و دنبال‌کردن مسیرهای چندگامی را فراهم می‌کند. همچنین، افزودن یا اصلاح یک رابطه بدون آموزش دوباره مدل زبانی امکان‌پذیر است.

در روش‌های بازیابی‌محور\footnote{\lr{Retrieval-based Methods}}، یک بازیاب بیرونی بخش مرتبط گراف یا اسناد متصل به آن را انتخاب می‌کند و نتیجه را در اختیار مدل زبانی قرار می‌دهد. مزیت این الگو آن است که فرایند بازیابی از تولید پاسخ جداست و می‌توان کیفیت زیرگراف، میزان پوشش شواهد و هزینه بازیابی را مستقل ارزیابی کرد. در مقابل، عملکرد روش به دقت بازیاب وابسته است؛ حذف یک گره یا مسیر مهم در این مرحله ممکن است با توان استدلالی مدل جبران نشود. \lr{G-Retriever}، \lr{HippoRAG 2} و در مقیاسی محدودتر \lr{LinearRAG} نمونه‌هایی از این جهت‌گیری‌اند که با سازوکارهای متفاوت، ارتباط ساختاری میان قطعه‌های متنی را در بازیابی وارد می‌کنند \cite{he2024gretriever,gutierrez2025hipporag2,zhuang2026linearrag}.

در روش‌های عامل‌محور\footnote{\lr{Agent-based Methods}}، مدل زبانی تنها مصرف‌کننده نتیجه بازیابی نیست، بلکه در انتخاب گام بعدی جست‌وجو نقش دارد. مدل از مفاهیم اولیه شروع می‌کند، رابطه‌ها و موجودیت‌های محتمل را ارزیابی می‌کند و این چرخه تا یافتن شواهد کافی یا رسیدن به حد توقف ادامه می‌یابد. روش \lr{ToG} این ایده را با پیمایش تکرارشونده مسیرهای گرافی به کار می‌گیرد و \lr{ToG-2} بازیابی گراف و متن را به صورت رفت‌وبرگشتی ترکیب می‌کند \cite{sun2024think,ma2025tog2}. این الگو برای پرسش‌های چندگامی انعطاف بیشتری دارد، اما به فراخوانی‌های متعدد مدل وابسته است و در نتیجه زمان پاسخ و هزینه آن افزایش می‌یابد. افزون بر این، انتخاب نادرست در یک گام می‌تواند ادامه پیمایش را به مسیر نامرتبط هدایت کند.

الگوی سوم، استدلال مقید به گراف\footnote{\lr{Graph-constrained Reasoning}} است. در این الگو، ساختار گراف تنها در ورودی مدل قرار نمی‌گیرد، بلکه گزینه‌های مجاز هنگام تولید مسیر استدلال را محدود می‌کند. روش \lr{GCR} مسیرهای گراف را در نمایه‌ای به نام \lr{KG-Trie} ذخیره می‌کند و مدل را وادار می‌سازد دنباله‌ای سازگار با مسیرهای موجود در گراف تولید کند \cite{luo2025gcr}. این قید از ساخت رابطه‌ای که در گراف وجود ندارد جلوگیری می‌کند و ردیابی مسیر را دقیق‌تر می‌سازد؛ با این حال، اگر دانش لازم در گراف ثبت نشده باشد، محدودیت ساختاری نمی‌تواند آن را جبران کند.

تفاوت این سه الگو در محل اعمال ساختار گراف است. در روش بازیابی‌محور، گراف بر انتخاب شواهد اثر می‌گذارد؛ در روش عامل‌محور، گراف فضای جست‌وجوی تعاملی مدل را شکل می‌دهد؛ و در روش مقید، مسیر تولیدشده باید با ساختار گراف سازگار باشد. در نتیجه، هر الگو تعادل متفاوتی میان پوشش شواهد، انعطاف استدلال، وفاداری مسیر و هزینه محاسباتی برقرار می‌کند. مقایسه منصفانه آن‌ها نیازمند داده، گراف و معیارهای ارزیابی مشترک است.

جهت‌گیری این پژوهش، بازپیاده‌سازی و انطباق نماینده‌هایی از این الگوها برای تشخیص توضیح‌پذیر سلامت روان است. همه روش‌ها بر وظایف منتخب \lr{IMHI} و یک گراف دانش دامنه‌ای مشترک ارزیابی می‌شوند تا اثر سازوکار ادغام از تفاوت داده و دانش جدا شود. مقایسه بر درستی برچسب، وفاداری و کیفیت توضیح، زمان پاسخ و هزینه متنی متمرکز است. افزون بر معیارهای معمول، ادعاها و روابط بیان‌شده در توضیح با چارچوب \lr{GraphEval} بررسی می‌شوند تا مشخص شود هر پاسخ تا چه اندازه به شواهد قابل ردیابی متکی است.

\section{جمع‌بندی}

در این فصل، مسئله تشخیص توضیح‌پذیر سلامت روان از داده‌های متنی و تفاوت آن با طبقه‌بندی صرف بیان شد. سپس نقش مدل‌های زبانی بزرگ در تولید هم‌زمان برچسب و توضیح، مسئله توهم اطلاعاتی و نیاز به سنجش مستقل وفاداری توضیح بررسی گردید. بازیابی افزوده به عنوان راهی برای اتصال مدل به دانش بیرونی معرفی شد و نشان داده شد که بازیابی متنی، به دلیل رتبه‌بندی مستقل قطعه‌ها و ناتوانی در نمایش صریح روابط، در پرسش‌های چندگامی با محدودیت روبه‌روست.

در ادامه، ساختار گرافی و سه الگوی بازیابی‌محور، عامل‌محور و استدلال مقید از یکدیگر تفکیک شدند. این الگوها از نظر نحوه انتخاب شواهد، میزان دخالت مدل در پیمایش، وفاداری مسیر و هزینه اجرایی تفاوت دارند و مبنای مقایسه این پژوهش را تشکیل می‌دهند. فصل بعد، پژوهش‌های مرتبط و روش‌های مبنا را بر پایه همین تقسیم‌بندی بررسی می‌کند و مزایا و محدودیت‌های هر خانواده را با جزئیات بیشتری توضیح می‌دهد.
